\chapter[Imagen aritmética finita y factorización \(468=18\cdot26\)]{Imagen aritmética finita y factorización \(468=18\cdot26\)}
\label{chap:alfabeto-42}
\index{código decimal admisible}
\index{dígito de control}

La imagen en base \(1000\) de la aplicación finita que asigna los bloques decimales es un
subconjunto propio de los mil bloques decimales. La ley de evolución definida en el
\cref{chap:tpk-definicion} y la aritmética de \(\mathbb Z/9\mathbb Z\)
determinan una imagen de 42 ternas, organizada en seis clases de siete. La
misma separación de los observables aditivo y multiplicativo factoriza después las 468
salidas en 18 firmas aditivas y 26 firmas multiplicativas.

\section{Primitivas y recurrencia del emisor finito}

\begin{definition}[Datos de partida de \(T_{\min}\)]
\label{pf84:E02-15}
Sea \(D=\{N,E,S,O\}\), con vectores cardinales
\[
 v(N)=(-1,0),\quad v(E)=(0,1),\quad
 v(S)=(1,0),\quad v(O)=(0,-1)
 \qquad\text{en }(\mathbb Z/9\mathbb Z)^2.
\]
La firma temporal y el índice multiplicativo se fijan por
\[
 \operatorname{sgn}_9(t)=
 \begin{cases}
 +1,&t\bmod9\in\{1,4,7\},\\
 -1,&t\bmod9\in\{2,5,8\},\\
 0,&t\bmod9\in\{0,3,6\},
 \end{cases}
\]
y
\[
 \operatorname{ind}_2(1,2,4,8,7,5)=(0,1,2,3,4,5),
\]
con valor cero fuera de las unidades. Una entrada de \(T_{\min}\) es un par
de semillas
\[
 I_9:=\{0,1,\ldots,8\},\qquad
 s_+=(i_+,j_+,d_+),\qquad s_\times=(i_\times,j_\times,d_\times)
 \in I_9^2\times D.
\]
Estas tablas, los lectores \(L_\Sigma,L_\Pi\), los seis bloques de nueve
instantes y la inversión de ambos vectores cuando \(27\mid t\) son todos los
datos de partida. No se introduce \(\pi\), \(e\), \(\varphi\) ni ningún otro
valor objetivo.
\end{definition}

\begin{algorithm}[Recurrencia de 54 pasos]
\label{pf84:E02-16}
Al comenzar, los dos cursores ocupan las posiciones de sus semillas y llevan
los vectores \(v(d_+)\), \(v(d_\times)\). Para
\(t=1,\ldots,54\), en este orden:
\begin{enumerate}
  \item si \(\operatorname{sgn}_9(t)=+1\), se acumula
  \(L_\Sigma(i_+(t),j_+(t))\) y se desplaza el cursor aditivo;
  \item si \(\operatorname{sgn}_9(t)=-1\), se acumula
  \(\operatorname{ind}_2(L_\Pi(i_\times(t),j_\times(t)))\) y se desplaza el
  cursor multiplicativo;
  \item si \(\operatorname{sgn}_9(t)=0\), ambos cursores permanecen quietos;
  \item después de esa actualización, ambos vectores se invierten si
  \(27\mid t\).
\end{enumerate}
Los acumuladores se reinician al comienzo de cada bloque
\(9(r-1)+1\leq t\leq9r\). Si \(S_r,E_r\) son sus valores finales, se emite
\[
 U_r=100\bar S_r+10(\bar S_r+\bar E_r)_{10}
       +(3S_r+5E_r+1)_{10},
 \qquad
 \bar S_r=S_r\bmod10,\quad \bar E_r=E_r\bmod10,
\]
y \(w_r=(-U_r)\bmod3\). Así se obtiene
\[
 T_{\min}(s_+,s_\times)=(U_1,\ldots,U_6)
 \in\mathcal U_6^{\rm cat},
 \qquad
 \Psi(U_6)=(-U_6)\bmod3=:w_6\in\mathbb F_3^6.
\]
\end{algorithm}

\begin{criterion}[Control dual independiente del valor objetivo]
\label{pf84:E02-17}
La recurrencia admite dos formulaciones independientes: factorización de
canales y simulación conjunta. El recorrido exhaustivo de las
\(104\,976\) entradas produce el mismo sexteto por ambas vías para cada
semilla, sin consultar una constante objetivo. Una discrepancia en una sola
semilla refutaría la equivalencia.
\end{criterion}

\section{Dinámica bivariada y observables por bloques}

Trabajamos en el toro \(\mathbb Z_9^2\). Con la convención precedente, una
semilla completa es
\[
\omega=(p_+,p_\times,h_+,h_\times)
\in(\mathbb Z_9^2)^2\times\{N,E,S,O\}^2,
\]
de modo que
\[
|\Omega|=9^4\cdot4^2=104\,976.
\]
En cada ventana de nueve pasos, los tiempos \(1,4,7\) evalúan el observable aditivo,
los tiempos \(2,5,8\) evalúan el observable multiplicativo y los tiempos \(3,6,9\)
incrementan el contador neutral. Si \(S_m,E_m,T_m\) son los acumuladores
de la ventana \(m\), el observable decimal es
\[
d_1\equiv S_m,\qquad
d_2\equiv S_m+E_m,\qquad
d_3\equiv3S_m+5E_m+7T_m\pmod{10},
\]
y la terna decimal es \(U_m=100d_1+10d_2+d_3\).

\begin{lemma}[Contador neutral]
Para toda ventana, \(T_m=3\).
\end{lemma}
\begin{proof}
Cada bloque de nueve tiempos contiene exactamente los tres tiempos neutrales
\(3,6,9\). Por tanto el contador se incrementa tres veces.
\end{proof}

\section{Imagen del acumulador aditivo: \(\operatorname{Im}S_m=\{6,9,12,15,18,21,24\}\)}

Sea \(s=i+j\pmod9\). El observable aditivo visible es
\[
f(s)=\operatorname{dr}_9(s+2)=(2,3,4,5,6,7,8,9,1).
\]
Una traslación cardinal modifica \(s\) en \(+1\) o en \(-1\). En una ventana se
leen tres términos consecutivos; invertir la orientación sólo invierte su
orden. Las nueve sumas cíclicas son
\[
9,12,15,18,21,24,18,12,6.
\]

\begin{lemma}[Imagen del acumulador aditivo]
\[
S_m\in\{6,9,12,15,18,21,24\}.
\]
En particular, \(d_1=S_m\bmod10\) recorre exactamente
\(\{1,2,4,5,6,8,9\}\).
\end{lemma}
\nopagebreak[4]
\begin{proof}
La lista anterior enumera todas las posiciones iniciales de la trayectoria y las
dos orientaciones. Sus residuos decimales son siete y distintos.\qedhere
\end{proof}

\section{Imagen del acumulador multiplicativo: \(\operatorname{Im}E_m=\{0,1,3,5,7,9\}\)}

El grupo de unidades \((\mathbb Z/9\mathbb Z)^\times\) es cíclico de orden
seis. Fijamos la función de coordenada multiplicativa
\[
\ell_\times(1,2,3,4,5,6,7,8,9)=(0,1,0,2,5,0,4,3,0).
\]
Si la fila o columna recorrida posee un factor divisible por tres, cada
evaluación pertenece a \(\{3,6,9\}\) y aporta cero. Si el factor fijo es una
unidad, las sumas de tres términos consecutivos de
\(b\mapsto\ell_\times(\operatorname{dr}_9(ab))\) son:
\[
\begin{array}{c|c}
a&\text{suma posible}\\
\midrule
1&\{1,3,7,9\}\\
2&\{3,5,9\}\\
4&\{1,5,7\}\\
5&\{1,5,7\}\\
7&\{3,5,9\}\\
8&\{1,3,7,9\}.
\end{array}
\]

\begin{lemma}[Imagen del acumulador multiplicativo]
\[
E_m\in\{0,1,3,5,7,9\}.
\]
\end{lemma}
\begin{proof}
El caso no invertible aporta \(0\). La tabla finita de las seis unidades
contiene todas las posibilidades del caso invertible y su unión es
\(\{1,3,5,7,9\}\).
\end{proof}

\section{Congruencia decimal del alfabeto ternario: unicidad del 3.er dígito}

\begin{theorem}[Congruencia de control]
Toda terna \(\overline{d_1d_2d_3}\) emitida por el sistema satisface
\[
\boxed{d_3\equiv5d_2-2d_1+1\pmod{10}.}
\]
\end{theorem}
\begin{proof}
De \(d_2\equiv S_m+E_m\) se obtiene
\(E_m\equiv d_2-d_1\pmod{10}\). Como \(T_m=3\),
\[
d_3\equiv3d_1+5(d_2-d_1)+21
\equiv5d_2-2d_1+1\pmod{10}.\qedhere
\]
\end{proof}

\begin{theorem}[Imagen exacta de la aplicación decimal]
El sistema produce exactamente \(7\cdot6=42\) ternas decimales. Agrupadas por
\(E=d_2-d_1\pmod{10}\), son:
\[
\begin{array}{c|rrrrrrr}
E&\multicolumn{7}{c}{\text{ternas decimales}}\\
\midrule
0&114&227&443&556&669&885&998\\
1&129&232&458&561&674&890&903\\
3&149&252&478&581&694&810&923\\
5&169&272&498&501&614&830&943\\
7&189&292&418&521&634&850&963\\
9&109&212&438&541&654&870&983.
\end{array}
\]
\end{theorem}
\begin{proof}
Hay siete posibilidades para \(d_1\) y seis para \(E\). Para cada par,
\(d_2\equiv d_1+E\) y la congruencia de control fija un único \(d_3\),
de modo que existen a lo sumo 42 ternas. La tabla aplica estas dos fórmulas
a los 42 pares y no contiene repeticiones. La enumeración de las
104\,976 semillas encuentra todas las ternas de la tabla y ninguna fuera de
ella; la cota se alcanza.
\end{proof}

\section{Factorización de los \(2\) observables}

Sea \(\mathcal S=\mathbb Z_9^2\times\{N,E,S,O\}\), con
\(|\mathcal S|=324\). A cada semilla aditiva \(P\in\mathcal S\) se le
asocia la firma de seis ventanas
\[
s(P)=(S_1,\ldots,S_6)\pmod{10},
\]
y a cada semilla multiplicativa \(T\in\mathcal S\),
\[
e(T)=(E_1,\ldots,E_6)\pmod{10}.
\]
Las dos firmas se combinan coordenada a coordenada mediante
\[
G(s,e)_k=
100s_k+10(s_k+e_k)_{10}+(3s_k+5e_k+1)_{10}.
\]
La aplicación \(G\) es inyectiva: desde una terna se recupera
\(s_k=d_1\) y \(e_k=d_2-d_1\pmod{10}\), mientras el tercer dígito verifica
la pertenencia a la imagen.

\begin{lemma}[Clasificación exacta del canal aditivo]
\label{pf84:E02-18}
\[
 |\operatorname{im}s|=18,
\]
y cada firma aditiva tiene exactamente \(18\) preimágenes.
\end{lemma}
\begin{proof}
La firma sólo depende de
\((i+j\bmod9,\varepsilon)\), donde \(\varepsilon=+1\) para \(E,S\) y
\(-1\) para \(N,O\). Los \(9\cdot2\) pares producen firmas distintas; cada
uno posee nueve posiciones en su diagonal y dos direcciones, es decir,
\(18\) preimágenes. Un número distinto de firmas o una clase de tamaño
distinto de \(18\) refutaría la clasificación.
\end{proof}

\begin{lemma}[Clasificación exacta del canal multiplicativo]
\label{pf84:E02-19}
\[
 |\operatorname{im}e|=26.
\]
Hay \(24\) clases de tamaño \(8\), una de tamaño \(24\) y una de tamaño
\(108\).
\end{lemma}
\begin{proof}
La enumeración de las \(9^2\cdot4=324\) semillas mediante la tabla explícita
de \(\ell_\times\) da \(26\) sextetos
\((E_1,\ldots,E_6)\) y el histograma
\[
 24\cdot8+24+108=324.
\]
Una clase adicional, ausente o de tamaño diferente refutaría el censo.
\end{proof}

\begin{theorem}[Factorización exacta e inyectiva del catálogo]
\label{pf84:E02-20}
\[
|\operatorname{im}s|=18,\qquad
|\operatorname{im}e|=26,\qquad
|\operatorname{im}G|=468=18\cdot26.
\]
Cada firma aditiva tiene 18 preimágenes. Las firmas multiplicativas poseen
el histograma
\[
24\text{ clases de tamaño }8,\qquad
1\text{ de tamaño }24,\qquad
1\text{ de tamaño }108.
\]
\end{theorem}
\begin{proof}
Los dos lemas precedentes dan los cardinales de las firmas. En cada componente
emitida, la cifra de centenas es \(\bar S_r\), mientras la de decenas es
\(\bar S_r+\bar E_r\bmod10\). Por tanto,
\[
 \bar S_r=\left\lfloor\frac{U_r}{100}\right\rfloor,\qquad
 \bar E_r=
 \left(\left\lfloor\frac{U_r}{10}\right\rfloor-\bar S_r\right)\bmod10.
\]
La cifra de unidades verifica \(3S_r+5E_r+1\bmod10\); no introduce otra
identificación. Así, \(G\) es inyectiva y su imagen tiene
\(18\cdot26=468\) elementos. Las \(24\) clases multiplicativas de tamaño
\(8\), combinadas con las \(18\) firmas aditivas, producen \(432\) sextetos
de multiplicidad \(144\); las otras dos clases producen \(18\) sextetos de
multiplicidad \(432\) y \(18\) de multiplicidad \(1944\). En particular,
\[
 432\cdot144+18\cdot432+18\cdot1944=104\,976.
\]
Dos pares de firmas con el mismo \(U_6\), o un histograma que no recupere el
dominio completo, refutarían el teorema.
\end{proof}

La composición completa factoriza, por tanto, como
\[
\mathcal S^2\xrightarrow{s\times e}
\operatorname{im}s\times\operatorname{im}e
\xrightarrow{G}\mathcal U_6
\xrightarrow{\rho_3}\mathcal W_6\subset\mathbb F_3^6.
\]
La última proyección posee 243 valores. El primer descenso conserva la
separación suma--producto; el segundo conserva 468 estados decimales; el
tercero elimina la coordenada de fibra y retiene una palabra ternaria.

\begin{corollary}[Censo de la proyección observable]
\label{cor:censo-proyeccion-observable}
\label{pf84:E02-21}
La enumeración completa determina la cadena de cardinales
\[
 104\,976\longrightarrow468\longrightarrow243\subset729.
\]
Los \(468\) elementos intermedios son sextetos de enteros comprendidos entre
cero y \(999\); sólo su imagen final pertenece a \(\mathbb F_3^6\). Las
multiplicidades de las fibras del primer mapa satisfacen
\[
 432\cdot144+18\cdot432+18\cdot1944=104\,976.
\]
\end{corollary}

\begin{proof}
El teorema de factorización da \(18\cdot26=468\), y la enumeración
exhaustiva de las \(104\,976\) semillas produce exactamente las tres
multiplicidades indicadas. La reducción componente a componente recorre las
\(243\) palabras registradas en el catálogo ternario. Los catálogos de
firmas y rutas adjuntos reproducen ambas enumeraciones.
\end{proof}

\begin{proposition}[Microfibra quíntuple de la palabra \(010211\)]
\label{pf84:E02-22}
Para la reducción restringida
\[
 \Psi:\mathcal U_6^{\rm cat}\longrightarrow\mathbb F_3^6,
 \qquad \Psi(U_6)=(-U_6)\bmod3,
\]
la fibra de \(010211\) consta exactamente de los cinco sextetos siguientes:
\begingroup
\scriptsize
\[
\begin{array}{@{}cclc@{}}
\toprule
U_6&\sigma_+&\sigma_\times&\text{multiplicidad}\\
\midrule
(501,614,498,169,272,272)&(5,6,4,1,2,2)&(5,5,5,5,5,5)&432\\
(810,923,870,169,272,272)&(8,9,8,1,2,2)&(3,3,9,5,5,5)&144\\
(810,983,810,169,272,272)&(8,9,8,1,2,2)&(3,9,3,5,5,5)&144\\
(870,923,810,169,272,272)&(8,9,8,1,2,2)&(9,3,3,5,5,5)&144\\
(870,923,810,418,674,521)&(8,9,8,4,6,5)&(9,3,3,7,1,7)&144\\
\bottomrule
\end{array}
\]
\endgroup
Son cinco preimágenes del catálogo \(U_6\), no cinco historias completas
TPK. Conservan firmas y multiplicidades distintas aunque compartan la misma
palabra visible.
\end{proposition}

\begin{proof}
El barrido exhaustivo del catálogo de \(468\) sextetos selecciona exactamente
las cinco filas de la tabla. La función de recuperación de firmas reproduce
los dos sextetos de cada fila y la suma de multiplicidades es
\(432+4\cdot144=1008\). La enumeración comprueba además que cada celda es el
producto cartesiano de sus semillas de canal. Un cardinal distinto de cinco
o la identificación de dos filas pese a sus datos enriquecidos refutaría el
enunciado.
\end{proof}
